A point path follows one chosen input through a sequence of maps.  A ripple
pencil follows an entire output condition backward through the same maps.  The
upper half-plane makes this dual reading elementary because M\"obius maps carry
generalized circles to generalized circles.

\subsection{Poles and standard horocycles}

Let
\begin{equation}\label{eq:p0-mobius-map}
  F(z)=\frac{Az+B}{Cz+D}
\end{equation}
be a real fractional linear map with
\[
  \Delta=AD-BC>0.
\]
The positivity convention ensures that $F$ maps the upper half-plane to itself.
The projective pole is
\[
  p_F=F^{-1}(\infty)
  =\begin{cases}
     -D/C,&C\ne0,\\
     \infty,&C=0.
   \end{cases}
\]
For $t>0$, denote the standard horocycle based at infinity by
\[
  \mathcal H_t=\{w\in\mathbb H:\im w=t\}.
\]

\begin{definition}[Ripple pencil]
\label{def:p0-ripple-pencil}
The \emph{ripple pencil} of $F$ is the one-parameter family
\[
  \mathcal R(F)
  =\{F^{-1}(\mathcal H_t):t>0\}.
\]
It is a family of horocycles based at the pole $p_F$.  The word ``ripple'' is a
visual name for this standard M\"obius pullback family.
\end{definition}

\begin{theorem}[Line-to-circle transition]
\label{thm:p0-ripple-pencil-formula}
Let $F$ be as in \eqref{eq:p0-mobius-map}.

\begin{enumerate}
  \item If $C=0$, every member of $\mathcal R(F)$ is a horizontal line; the
  pencil is based at $\infty$.

  \item If $C\ne0$, the member with parameter $t$ is the Euclidean circle
  \begin{equation}\label{eq:p0-ripple-circle}
    \left(x+\frac DC\right)^2
    +\left(y-\frac{\Delta}{2tC^2}\right)^2
    =\left(\frac{\Delta}{2tC^2}\right)^2.
  \end{equation}
  All these circles are tangent to the real boundary at the common point
  $p_F=-D/C$ and are nested as $t$ varies.
\end{enumerate}
\end{theorem}

\begin{proof}
For $z=x+iy$, direct calculation gives the standard identity
\begin{equation}\label{eq:p0-imaginary-mobius}
  \im F(z)=\frac{\Delta y}{|Cz+D|^2}.
\end{equation}
The inverse image of $\mathcal H_t$ is therefore
\[
  \Delta y=t\bigl((Cx+D)^2+C^2y^2\bigr).
\]
If $C=0$, this reduces to a horizontal line.  If $C\ne0$, divide by
$tC^2$ and complete the square in $y$:
\[
  \left(x+\frac DC\right)^2+y^2
  -\frac{\Delta}{tC^2}y=0,
\]
which is exactly \eqref{eq:p0-ripple-circle}.  The center lies vertically
above $-D/C$ at a height equal to the radius, so the circle is tangent to the
boundary there.  The radius is proportional to $t^{-1}$, proving nesting.
\end{proof}

The theorem gives a geometric interpretation of the lower-left matrix entry.
The affine condition $C=0$ places the pole at infinity and degenerates the
ripple circles to parallel lines.  A projective step with $C\ne0$ moves the pole
to a finite boundary point and curves the same horocyclic family into visible
ripples.

\begin{example}[Arithmetic inversion]
\label{ex:p0-inversion-ripples}
For
\[
  J(z)=-\frac1z,
  \qquad
  J\leftrightarrow
  \begin{pmatrix}0&-1\\1&0\end{pmatrix},
\]
one has $p_J=0$ and
\begin{equation}\label{eq:p0-inversion-circles}
  x^2+\left(y-\frac1{2t}\right)^2
  =\left(\frac1{2t}\right)^2.
\end{equation}
Thus inversion turns the horizontal horocycles based at infinity into a nested
family of circles tangent at zero.  The map exchanges the two ideal points
$0$ and $\infty$.
\end{example}

\begin{figure}[t]
\centering
\begin{tikzpicture}[
  x=1.05cm,y=1.05cm,
  axis/.style={-{Latex[length=2mm]},thick,black!70},
  ripple/.style={thick,blue!60!black},
  geod/.style={dashed,gray!55},
  every node/.style={font=\small}
]
  \fill[blue!2] (-4.2,0) rectangle (4.2,4.0);
  \draw[axis] (-4.35,0)--(4.45,0) node[right] {$x$};
  \draw[axis] (0,-0.08)--(0,4.2) node[above] {$y$};
  \foreach \r in {0.45,0.75,1.15,1.7}
    {\draw[ripple] (0,\r) circle (\r);}
  \draw[geod] (0,0)--(0,3.75);
  \fill[orange!80!black] (0,0) circle (2.2pt)
    node[below right] {pole $p_J=0$};
  \node[blue!55!black,anchor=west] at (1.85,2.75)
    {$J^{-1}(\mathcal H_t)$};
  \node[anchor=west,font=\scriptsize] at (1.85,2.3)
    {larger circles correspond to smaller $t$};
\end{tikzpicture}
\caption{The ripple pencil of inversion.  Every circle is a horocycle tangent
at the projective pole $0$; the affine horizontal-line pencil is recovered
when the pole is moved to infinity.}
\label{fig:p0-inversion-ripple-pencil}
\end{figure}

\subsection{Zero--pole pencils in the complex chart}

There is a second elementary pencil attached to the same operator.  Suppose
$A,C\ne0$ and write its zero and pole as
\[
  q_F=-\frac BA,
  \qquad
  p_F=-\frac DC.
\]
Then the output level set $|F(z)|=r$ is equivalent to
\begin{equation}\label{eq:p0-apollonius-pencil}
  |z-q_F|
  =r\left|\frac CA\right|\,|z-p_F|.
\end{equation}
These are Apollonius circles for the ordered zero--pole pair.  The horocyclic
ripple pencil and the Apollonius pencil encode different output observables,
$\im F$ and $|F|$, but both show that a projective one-hole operator naturally
carries more than a point orbit: it carries a zero, a pole, and a family of
level curves organized by them.

\subsection{Ripple histories of right-expanded expressions}

Let a right-expanded expression determine nondegenerate real M\"obius sections
$F_1,\ldots,F_n$, in chronological order, and write
\[
  G_k=F_k\circ\cdots\circ F_1,
  \qquad
  G_k(z)=\frac{A_kz+B_k}{C_kz+D_k}.
\]
At stage $k$, the arithmetic value of a chosen input follows the point orbit
$G_k(z_0)$.  Simultaneously, the output horocycles pull back to
\[
  \mathcal R_k=\mathcal R(G_k).
\]
The tangency point
\[
  p_k=G_k^{-1}(\infty)
  =-\frac{D_k}{C_k}
\]
when $C_k\ne0$ is the current pole of the prefix operator.

Thus a general right-expansion history should not be pictured as one fixed
family of concentric circles.  Each prefix has its own ripple pencil, and the
common tangency point of that pencil may move with $k$.  The appropriate
geometric state is
\[
  \bigl(G_k,\,G_k(z_0),\,p_k,\,\mathcal R_k\bigr):
\]
the operator, the selected point, the pole, and the pulled-back wavefront
family.

\begin{proposition}[Covariant--contravariant duality]
\label{prop:p0-path-ripple-duality}
For every prefix operator $G_k$ and every output horocycle $\mathcal H_t$,
\[
  z_0\in G_k^{-1}(\mathcal H_t)
  \quad\Longleftrightarrow\quad
  G_k(z_0)\in\mathcal H_t.
\]
Hence the point path and the ripple history are two readings of the same
operator word: one pushes the input forward, the other pulls the output
condition backward.
\end{proposition}

\begin{proof}
This is the defining property of inverse image.
\end{proof}

The elementary nature of the proposition is precisely its value.  No new
geometry is inserted between the two pictures.  Their unity is already present
in the action of one M\"obius matrix.
